% GENERATED FILE. Edit canonical lessons, then run npm run generate:books.
% canonical-source-hash: fnv1a64:cf0cbdbd9554a476
% canonical-lessons: FA-C318-jenayat, FA-C318-qatl, FA-C318-qatel, FA-C318-hamle-kardan, FA-C318-aslahe

\chapter{Serious crime}
\label{ch:fa-serious-crime}
\hlchaptermodality{318}

\begin{chapteropening}
\textbf{By the end of this chapter:} \emph{I can talk about serious crimes: a crime, a murder, a murderer, an attack and a weapon.}

Complete the last lesson of chapter 318: I can talk about serious crimes: a crime, a murder, a murderer, an attack and a weapon.
\end{chapteropening}

\section[\texorpdfstring{\fa{جنایت} (jenâyat)}{جنایت (jenâyat)}]{\fa{جنایت} (jenâyat) — a serious crime, a felony}

\label{lesson:FA-C318-jenayat}



\begin{quote}
Before the new one: say the Persian for to disappear, then the Persian for to be found out.
\end{quote}

\subsection*{You'll want to know: \fa{جنایت}}
\textbf{\fa{جنایت}} (\emph{jenâyat}) — \textquotedblleft{}a serious crime, a felony\textquotedblright{}.

\begin{quote}
\fa{این} \fa{جنایت} \fa{همه} \fa{شهر} \fa{را} \fa{شوکه} \fa{کرد}.
\end{quote}

\begin{quote}
\emph{In jenâyat hame-ye shahr râ shoke kard.}
\end{quote}

\begin{quote}
This crime shocked the whole city.
\end{quote}

For violent crimes. \emph{jenâyatkâr} \textquotedblleft{}a criminal\textquotedblright{}.

\begin{cousinweb}
\emph{jenâyat} is borrowed from Arabic \emph{jināya}, \textquotedblleft{}crime\textquotedblright{} (EtymDB).
\end{cousinweb}

\begin{grammarlens}[title={What you've built}]
One more word for telling, explaining and describing at B1.
\end{grammarlens}

\subsection*{Your turn}
\begin{itemize}
\raggedright
  \item \emph{Say it:} \emph{jenâyat}
  \item \emph{Say it:} \emph{jenâyat}, once more
  \item \emph{Say it:} a sentence of your own with \emph{jenâyat}
  \item \emph{Recall:} say the Persian for to disappear, then the Persian for to be found out, then say \emph{jenâyat} again
\end{itemize}

\begin{tcolorbox}[breakable,colback=teal!4,colframe=teal!35!black,title={Before you move on}]
What does \fa{جنایت} mean? (\textquotedblleft{}A serious crime, a felony\textquotedblright{}.) Say it once more.
\end{tcolorbox}
\section[\texorpdfstring{\fa{قتل} (qatl)}{قتل (qatl)}]{\fa{قتل} (qatl) — a murder, a killing}

\label{lesson:FA-C318-qatl}



\begin{quote}
Before the new one: say the Persian for to be found out, then the Persian for a serious crime.
\end{quote}

\subsection*{You'll want to know: \fa{قتل}}
\textbf{\fa{قتل}} (\emph{qatl}) — \textquotedblleft{}a murder, a killing\textquotedblright{}.

\begin{quote}
\fa{قتل} \fa{در} \fa{یک} \fa{خانه} \fa{قدیمی} \fa{اتفاق} \fa{افتاد}.
\end{quote}

\begin{quote}
\emph{Qatl dar yek khâne-ye qadimi ettefâq oftâd.}
\end{quote}

\begin{quote}
The murder happened in an old house.
\end{quote}

Formal and legal. \emph{Be qatl residan} \textquotedblleft{}to be murdered\textquotedblright{}.

\begin{cousinweb}
\emph{qatl} is borrowed from Arabic \emph{qatl}, \textquotedblleft{}killing\textquotedblright{} (EtymDB).
\end{cousinweb}

\begin{grammarlens}[title={What you've built}]
One more word for telling, explaining and describing at B1.
\end{grammarlens}

\subsection*{Your turn}
\begin{itemize}
\raggedright
  \item \emph{Say it:} \emph{qatl}
  \item \emph{Say it:} \emph{qatl}, once more
  \item \emph{Say it:} a sentence of your own with \emph{qatl}
  \item \emph{Recall:} say the Persian for to be found out, then the Persian for a serious crime, then say \emph{qatl} again
\end{itemize}

\begin{tcolorbox}[breakable,colback=teal!4,colframe=teal!35!black,title={Before you move on}]
What does \fa{قتل} mean? (\textquotedblleft{}A murder, a killing\textquotedblright{}.) Say it once more.
\end{tcolorbox}
\section[\texorpdfstring{\fa{قاتل} (qâtel)}{قاتل (qâtel)}]{\fa{قاتل} (qâtel) — a murderer, a killer}

\label{lesson:FA-C318-qatel}



\begin{quote}
Before the new one: say the Persian for a serious crime, then the Persian for a murder.
\end{quote}

\subsection*{You'll want to know: \fa{قاتل}}
\textbf{\fa{قاتل}} (\emph{qâtel}) — \textquotedblleft{}a murderer, a killer\textquotedblright{}.

\begin{quote}
\fa{پلیس} \fa{هنوز} \fa{قاتل} \fa{را} \fa{پیدا} \fa{نکرده} \fa{است}.
\end{quote}

\begin{quote}
\emph{Polis hanuz qâtel râ peydâ nakarde ast.}
\end{quote}

\begin{quote}
The police have not found the murderer yet.
\end{quote}

Same Arabic root as \emph{qatl}: \emph{qâtel} is \textquotedblleft{}the one who kills\textquotedblright{}. Plural \emph{qâtelân}.

\begin{cousinweb}
\emph{qâtel} is borrowed from Arabic \emph{qātil}, \textquotedblleft{}killer\textquotedblright{} (EtymDB), the active participle of the same root as \emph{qatl}.
\end{cousinweb}

\begin{grammarlens}[title={What you've built}]
One more word for telling, explaining and describing at B1.
\end{grammarlens}

\subsection*{Your turn}
\begin{itemize}
\raggedright
  \item \emph{Say it:} \emph{qâtel}
  \item \emph{Say it:} \emph{qâtel}, once more
  \item \emph{Say it:} a sentence of your own with \emph{qâtel}
  \item \emph{Recall:} say the Persian for a serious crime, then the Persian for a murder, then say \emph{qâtel} again
\end{itemize}

\begin{tcolorbox}[breakable,colback=teal!4,colframe=teal!35!black,title={Before you move on}]
What does \fa{قاتل} mean? (\textquotedblleft{}A murderer, a killer\textquotedblright{}.) Say it once more.
\end{tcolorbox}
\section[\texorpdfstring{\fa{حمله} \fa{کردن} (hamle kardan)}{حمله کردن (hamle kardan)}]{\fa{حمله} \fa{کردن} (hamle kardan) — to attack}

\label{lesson:FA-C318-hamle-kardan}



\begin{quote}
Before the new one: say the Persian for a murder, then the Persian for a murderer.
\end{quote}

\subsection*{You'll want to know: \fa{حمله} \fa{کردن}}
\textbf{\fa{حمله} \fa{کردن}} (\emph{hamle kardan}) — \textquotedblleft{}to attack\textquotedblright{}.

\begin{quote}
\fa{یک} \fa{سگ} \fa{بزرگ} \fa{به} \fa{پسر} \fa{حمله} \fa{کرد}.
\end{quote}

\begin{quote}
\emph{Yek sag-e bozorg be pesar hamle kard.}
\end{quote}

\begin{quote}
A big dog attacked the boy.
\end{quote}

The one attacked takes \emph{be}. Present \emph{hamle mikonad}. \emph{Hamle-ye qalbi} \textquotedblleft{}a heart attack\textquotedblright{}.

\begin{cousinweb}
\emph{hamle} is borrowed from Arabic \emph{ḥamla} (EtymDB lists the Arabic word), usually glossed \textquotedblleft{}a charge, an attack\textquotedblright{}. \emph{kardan} \textquotedblleft{}to do\textquotedblright{} is inherited Persian.
\end{cousinweb}

\begin{grammarlens}[title={What you've built}]
One more word for telling, explaining and describing at B1.
\end{grammarlens}

\subsection*{Your turn}
\begin{itemize}
\raggedright
  \item \emph{Say it:} \emph{hamle kardan}
  \item \emph{Say it:} \emph{hamle kardan}, once more
  \item \emph{Say it:} a sentence of your own with \emph{hamle kardan}
  \item \emph{Recall:} say the Persian for a murder, then the Persian for a murderer, then say \emph{hamle kardan} again
\end{itemize}

\begin{tcolorbox}[breakable,colback=teal!4,colframe=teal!35!black,title={Before you move on}]
What does \fa{حمله} \fa{کردن} mean? (\textquotedblleft{}To attack\textquotedblright{}.) Say it once more.
\end{tcolorbox}
\section[\texorpdfstring{\fa{اسلحه} (aslahe)}{اسلحه (aslahe)}]{\fa{اسلحه} (aslahe) — a weapon, a gun}

\label{lesson:FA-C318-aslahe}



\begin{quote}
Before the new one: say the Persian for a murderer, then the Persian for to attack.
\end{quote}

\subsection*{You'll want to know: \fa{اسلحه}}
\textbf{\fa{اسلحه}} (\emph{aslahe}) — \textquotedblleft{}a weapon, a gun\textquotedblright{}.

\begin{quote}
\fa{مرد} \fa{اسلحه} \fa{نداشت}.
\end{quote}

\begin{quote}
\emph{Mard aslahe nadâsht.}
\end{quote}

\begin{quote}
The man did not have a weapon.
\end{quote}

An Arabic plural in origin, used as a singular in Persian. \emph{Aslahe-ye sard} is a knife or blade.

\begin{cousinweb}
\emph{aslahe} is borrowed from Arabic \emph{asliḥa} (EtymDB), the plural of \emph{silāḥ} \textquotedblleft{}weapon\textquotedblright{}; Persian treats it as a singular.
\end{cousinweb}

\begin{grammarlens}[title={What you've built}]
One more word for telling, explaining and describing at B1.
\end{grammarlens}

\subsection*{Your turn}
\begin{itemize}
\raggedright
  \item \emph{Say it:} \emph{aslahe}
  \item \emph{Say it:} \emph{aslahe}, once more
  \item \emph{Say it:} a sentence of your own with \emph{aslahe}
  \item \emph{Recall:} say the Persian for a murderer, then the Persian for to attack, then say \emph{aslahe} again
\end{itemize}

\begin{tcolorbox}[breakable,colback=teal!4,colframe=teal!35!black,title={Before you move on}]
What does \fa{اسلحه} mean? (\textquotedblleft{}A weapon, a gun\textquotedblright{}.) Say it once more.
\end{tcolorbox}
